\subsection{Why a ramification-nine base change of the full-nine examples does not solve the next gap}

\begin{proposition}\label{prop:B-RD25-basechange-obstruction}
\status{Proved; independent complete proof audit and exact certificate rerun by the parent passed.}
Let $K$ have absolute ramification index $24$, residue characteristic $3$, and perfect residue field.  If an additive elliptic curve $E/K$ has full rational $9$-torsion, then no extension of $K$ of ramification index $9$ makes its full $27$-torsion rational.  More generally, over its minimal good-reduction field after completed unramified base change, the next full division field has ramification degree at least $27$.

This rules out the proposed construction of an $(e,n)=(216,3)$ example by a ramification-nine base change of an $e=24$ full-nine example.  It does not rule out a new curve defined directly over a field of absolute index $216$; such a curve is constructed in Theorem~\ref{thm:B-RD26-wild}.
\end{proposition}
\begin{proof}
Pass first to completed unramified bases.  The proved tame criterion and the wild bound of Proposition~\ref{prop:B-RD20} force the semistability defect to be $12$.  Let $L/K$ be its minimal good-reduction field, so $v_L(3)=288$.  The good special fibre is supersingular.  Write $t$ for an integral formal parameter of the good model.

The universal valuation identity (B.12.4), together with the wild divisibility of the maximum valuation, gives the unique primitive-nine extrema
\[
                            A=21,\qquad B=1.
\]
Explicitly the only possible region is $r=1$, and
$288=6(A+27B)$ with $A\geq9B$, $3\mid A$, $B\geq1$
gives $(A,B)=(21,1)$; region $r=0$ would require
$A+3B=16$ with $3\mid A$, and is impossible.
Choose a basis $P,Q$ of $E[9]$ with $v_L(t(P))=21$ and $v_L(t(Q))=1$.

Formula (B.12.3) gives the minimum nonzero three-torsion valuation $9$.
The points of larger valuation form a line $C\subset E[3]$ of order three: two independent such points would make every nonzero three-torsion point have larger valuation.  The constant-term product for $[3](T)/T$ gives their common valuation $c$:
\[
                         288=6\cdot9+2c,\qquad c=117.
\]
All these torsion points are $L$-rational.

Let $R$ be a point with $3R=Q$, and let $S$ be any point with $3S=P$.
The Weierstrass factorization of $[3](T)$ is a unit times the product of $T-t(U)$ over $U\in E[3]$.  A point of valuation at least $9$ has image under $[3]$ of valuation at least $81$.  A point of smaller positive valuation $v$ has image of valuation $9v$.  Consequently
\[
                      v_L(t(R))=\frac19,\qquad
                      v_L(t(S))=\frac73.
\]
The field $F=L(R)$ has degree at most $9$, because $E[3]$ is rational; its ramification degree is at least $9$ by the valuation of $t(R)$.  Hence it is a totally ramified Galois extension of degree $9$.
The map $\sigma\mapsto\sigma R-R$ identifies
$\Gamma=\operatorname{Gal}(F/L)$ with $E[3]\simeq(\mathbf Z/3)^2$.
Put $z=t(R)$; it is a uniformizer of $F$.
Integral formal addition gives
\[
 v_F(\sigma z-z)=
 \begin{cases}
 1053,&0\ne\sigma R-R\in C,\\
 81,&\sigma R-R\notin C.
 \end{cases}
\]
Here the valuation of a formal difference equals that of the translating torsion point: the coefficient of that point in the difference is a unit.

Let $H\subset\Gamma$ correspond to $C$.  For every $0\ne\sigma\in H$ the lower ramification number is $1052$.  The elementary ramification estimate proved in Lemma~\ref{lem:B-RD20-wild} says
\[
 v_F(\sigma x-x)\geq v_F(x)+1052\qquad(x\in\mathfrak m_F).
\]
If $S$ belonged to $E(F)$, then $v_F(t(S))=21$; a nonzero translation of $S$ by three-torsion would have difference valuation either $81$ or $1053$, both less than $21+1052=1073$.  Thus every element of $H$ fixes $S$.  Set $F_1=F^H$; then $S$ would be $F_1$-rational, and $v_{F_1}(t(S))=7$.

We compute the ramification number of $F_1/L$ directly, without an appeal to a normalization of the upper numbering.  Its uniformizer
\[
                             w=N_{F/F_1}(z)
\]
satisfies $w=z^3+\epsilon$ with $v_F(\epsilon)\geq1055$.
For $\tau\notin H$, put $\delta=\tau z-z$, so $v_F(\delta)=81$.
Since $v_F(3)=2592$,
\[
 v_F\bigl((z+\delta)^3-z^3\bigr)
   =\min\{2592+2+81,\ 2592+1+162,\ 243\}=243.
\]
The correction $\tau\epsilon-\epsilon$ has valuation at least $1055$ and cannot cancel this term.  Hence
$v_{F_1}(\tau w-w)=81$, so the nontrivial automorphisms of $F_1/L$ have lower ramification number $80$.

For an element of valuation prime to $3$, the elementary ramification estimate is an equality: expand in the uniformizer with Teichm\"uller coefficients; the first term has coefficient equal to its exponent modulo $3$, which is nonzero.  Applied to $t(S)$ of valuation $7$, it gives
\[
                         v_{F_1}(\tau t(S)-t(S))=7+80=87.
\]
But $\tau S-S$ is a nonzero three-torsion point, so its formal coordinate has valuation $3\cdot9=27$ or $3\cdot117=351$ in $F_1$.  Neither equals $87$, a contradiction.
It follows that $F$ does not contain full $27$-torsion.  The full division field over $L$ is a $3$-extension, since the action is trivial modulo $9$; it contains $F$, and its degree is therefore at least $27$.

Finally, an extension $K'/K$ of ramification index $9$, followed by completed unramified base change, has degree $9$.  The compositum with $L$ has degree at most $9$ over $L$.  If it contained all $27$-torsion it would contradict the preceding argument.  This proves the asserted obstruction.
\end{proof}

\begin{remark}[Finite arithmetic certificate]
The exact certificate \path{code/verify_level27_basechange_obstruction.py}
enumerates every feasible $(A,B,r)$, checks the three-torsion and
next-division valuations, and verifies both ramification contradictions.
Its original run returned PASS in $0.006479513998783659$ seconds;
the parent independently reran it with PASS in $0.033169171001645736$ seconds.
These checks supplement the formal-group and field arguments above.
\end{remark}
